\section{正方形的渐伸线}
\markboth{46. 正方形的渐伸线}{46. 正方形的渐伸线}

蚊香盒的底面是一个正方形，拿一张旧报纸放在桌上，把蚊香盒放在报纸上，盒底的轮廓就在报纸上显示出一个正方形，记为 $ABCD$，并设其边长为 $a$。再取一根细线绕到蚊香盒外侧的下部，线的一端结上一支铅笔，使笔尖正好位于正方形的一个顶点 $A$ 处。然后手执铅笔拉紧细线，并且把线从盒底上逐渐松开，铅笔就在报纸上画出一条螺旋形曲线，这条曲线叫做正方形的渐伸线（图 46.1）。

\begin{wrapfigure}[11]{l}{0.43\textwidth}
  \centering
  \vspace{-4pt}
  \begin{tikzpicture}[scale=0.90, line cap=round, line join=round]
    \def\a{0.38}
    \coordinate (D) at (0, 0);
    \coordinate (C) at (\a, 0);
    \coordinate (B) at (\a, \a);
    \coordinate (A) at (0, \a);
    \coordinate (E) at (\a, {2*\a});
    \coordinate (F) at ({3*\a}, 0);
    \coordinate (G) at (0, {-3*\a});
    \coordinate (H) at ({-4*\a}, \a);

    % 正方形 ABCD
    \draw[thick] (D) -- (C) -- (B) -- (A) -- cycle;

    % 虚线辅助线
    \draw[densely dashed] (B) -- (E);
    \draw[densely dashed] (C) -- (F);
    \draw[densely dashed] (D) -- (G);
    \draw[densely dashed] (A) -- (H);

    % 渐伸线四段圆弧
    \draw[very thick, curveblue]
      (A) arc (180:90:\a)
          arc (90:0:{2*\a})
          arc (0:-90:{3*\a})
          arc (-90:-180:{4*\a});

    % 顶点与字母标签
    \fill (A) circle (1.2pt) node[above left=-1pt, font=\scriptsize] {$A$};
    \fill (B) circle (1.2pt) node[right, font=\scriptsize] {$B$};
    \fill (C) circle (1.2pt) node[below right=-1pt, font=\scriptsize] {$C$};
    \fill (D) circle (1.2pt) node[left, font=\scriptsize] {$D$};
    \node[above, font=\scriptsize] at (E) {$E$};
    \node[right, font=\scriptsize] at (F) {$F$};
    \node[below, font=\scriptsize] at (G) {$G$};
    \node[left, font=\scriptsize] at (H) {$H$};
  \end{tikzpicture}
  \caption*{图\ 46.1}
\end{wrapfigure}

根据作法知道，正方形的渐伸线是由一些不同半径的圆弧光滑连接而成的。例如在图 46.1 中，$\arc{AE}$ 的圆心是 $B$，半径是 $a$；$\arc{EF}$ 的圆心是 $C$，半径是 $2a$；$\arc{FG}$ 的圆心是 $D$，半径是 $3a$，等等。

还可用类似的方法画出圈数更多的螺旋形曲线，而且通过适当选择各个圆心的位置，螺旋形曲线的形状也可以千变万化。例如图 46.2 中画出的螺旋形曲线，也是由一些圆弧光滑地连接而成的，各段圆弧的圆心顺次是 $A$、$B$、$C$、$D$ 等等。这种画螺旋形曲线的方法非常简便，在机械制图和建筑工程中常被采用。

\begin{figure}[htbp]
  \centering
  \begin{tikzpicture}[scale=1.12, line cap=round, line join=round]
    % 宽阔平放的饱满蜗壳形四段圆弧螺旋，圆心顺次为 A, B, C, D
    \def\hB{0.72}
    \def\w{0.38}
    \def\hD{0.38}
    \pgfmathsetmacro{\rone}{2*\hB + \w}
    \pgfmathsetmacro{\rtwo}{\rone - \hB}
    \pgfmathsetmacro{\rthree}{\rtwo - \w}
    \pgfmathsetmacro{\rfour}{\hD}

    \coordinate (A) at (0, 0);
    \coordinate (B) at (0, \hB);
    \coordinate (C) at (\w, \hB);
    \coordinate (D) at (\w, \hD);

    % 水平底边及各段圆弧衔接半径辅助线
    \draw[thin] (-\rone, 0) -- (\w, 0);
    \draw[densely dashed, black!65] (A) -- (0, \rone);
    \draw[densely dashed, black!65] (B) -- (\rtwo, \hB);
    \draw[densely dashed, black!65] (C) -- (\w, 0);
    \draw[densely dashed, black!65] (D) -- (0, \hD);

    % 顺次以 A, B, C, D 为圆心光滑连接的四段 90° 圆弧螺旋
    \draw[very thick, curveblue]
      (-\rone, 0) arc (180:90:\rone)
                  arc (90:0:\rtwo)
                  arc (0:-90:\rthree)
                  arc (-90:-180:\rfour);

    % 圆心标记（全部置于清晰空白处，不压任何虚线与曲线）
    \fill (A) circle (1.1pt) node[below=1.5pt, font=\scriptsize] {$A$};
    \fill (B) circle (1.1pt) node[left=1.5pt, font=\scriptsize] {$B$};
    \fill (C) circle (1.1pt) node[above=1.5pt, font=\scriptsize] {$C$};
    \fill (D) circle (1.1pt) node[right=1.5pt, font=\scriptsize] {$D$};
  \end{tikzpicture}
  \caption*{图\ 46.2}
\end{figure}
